\section{Primera Entrevista}
\normalsize{ \indent
Con este proyecto, se desea controlar el legajo de
docentes y alumnos, así como que ellos puedan generar
formularios; entre otros, tener seguimiento de los
alumnos graduados del colegio. Para los legajos, podría
haber un seguimiento de cada docente/alumno con una
vista tipo tabla; inclusive, sería mejor tener los
archivos a mano junto con la tabla.
}
\newline
\normalsize{ \indent
Toda la información obligatoria de cada uno se halla en
los legajos, aunque en el futuro se podría necesitar
información extra. Por cierto, si se debe informar sobre
el estado de los legajos, sería únicamente a los docentes
por el correo que les otorga el colegio; se prefiere
informar más personalmente a los padres de los alumnos,
tal como teléfono fijo o celular.
}
\newline
\normalsize{ \indent
No hace falta nombrar todos los formularios de los
docentes y los alumnos, se cree que es suficiente con un
modelo digital de cada uno. Lo que si hace falta mencionar,
es que de los exalumnos se quiere conservar información
sobre su graduación de contacto.
}
\newline
\normalsize{ \indent
Los formularios se rellenan en la actualidad a mano,
aunque los modelos de los mismos se hallan en digital.
Cada vez que se necesite uno, se imprime el modelo y se
completa a mano; para después ser firmado, en caso de ser
necesario.
}
\newline
\normalsize{ \indent
En el sistema, se pretende que los docentes y alumnos
sólo puedan controlar sus legajos, perfiles y formularios.
Para el personal de Coordinación Escolar, se pretende
poder dar seguimiento de todos los legajos y formularios.
Por otra parte, la rectora se encarga de gestionar a los
docentes; designando los cargos y haciendo los relevos
necesarios.
}
\newline
\normalsize{ \indent
Del legajo, se puede digitalizar todas las partes, a
excepción del certificado médico de los alumnos; el cual
debe completar el médico que corresponda a cada uno, para
así guardar una copia digitalizada del mismo. De ahí, se
pueden guardar la información en una base de datos;
evitando que se deba reescribir una y otra vez la
información.
}
\newline
\normalsize{ \indent
Si se implementara algún sistema de comunicación, sería
de notificaciones; ya que existe la plataforma Acadeu,
la cual otorga un medio de comunicación entre los distintos
miembros del colegio. Además, no se pretende el envío de
mensajes entre los usuarios; el enfoque es sobre la gestión
de todo el papeleo en un solo lugar.
}
\section{Segunda Entrevista}
\normalsize{ \indent
Según su salario, los docentes se clasifican en jornales y
cargos; donde ambos roles pueden ser de un mismo docente.
De ahí, la cantidad de horas máxima que puede trabajar un
docente es 42 horas; incluyendo las horas trabajadas fuera
del colegio. Como la rectora es quién maneja los cargos,
ella también autoriza a los mismos; sin importar el título
que tengan en sí, siempre y cuando haya uno superior.
}
\newline
\normalsize{ \indent
Los tipos de relevos docentes que existen son de suplencia
y cambio de titular. La diferencia entre cada uno es que la
suplencia ocurre cuando el titular se halla de licencia,
mientras que el cambio del titular se refiere a una renuncia
de horas de un titular. Si un docente nuevo que releva a otro
por renuncia a las horas, puede que necesite un cambio de
horarios.
}
\newline
\normalsize{ \indent
El manejo de los horarios depende de Coordinación Escolar.
Dentro del año, pueden rotar todos los horarios; excepto
Módulo, el cual sólo tiene una rotación entre alumnos interna
(es decir, no afecta el horario). Para el cambio de horarios,
no se puede tener más de tres horas seguidas de una materia
pedagógica (se excluye Módulo, el cual si tiene más de tres
horas cátedra seguidas); además, el docente a cargo no debe
tener horarios externos que le impidan ejercer el cargo.
}
\newline
\normalsize{ \indent
En la actualidad, los cambios de horario se hacen con una
gestión completa de Coordinación Escolar; por lo cual, deben
revisar todas las condiciones para el cambio de horarios.
En otras palabras, por ahora es necesario revisar todos los
documentos manuscritos para poder controlar que se pueda
cambiar determinados horarios.
}
\newline
\normalsize{ \indent
Considerando todo lo dicho anteriormente, se tiene la mayoría
de razones de por qué se debe tener una validación del sistema
para el cambio de horarios. Lo único que amerita mencionar, es
que se prefiere automatizar completamente la selección de horarios
docentes; ya que puede alivianarse la carga que poseen los
trabajadores de Coordinación Escolar.
}
\newline
\normalsize{ \indent
El método de calificación de los alumnos es del 1 al 10, donde 6 es
la calificación para aprobar. Para el pase del año, los alumnos sólo
pueden tener hasta 3 materias desaprobadas por año. Por trimestre, debe
haber un mínimo de 3 calificaciones; en caso que una materia tenga menor
a ese, deberá contar con esa misma cantidad de todas formas.
}
\newline
\normalsize{ \indent
Los preceptores son los encargados de anotar la asistencia manualmente,
tanto de alumnos como docentes. Esta asistencia debe estar cubierta por
cada hora de cursado que haya, aunque cada presente dado abarca las horas
continuas que tenga una materia.
}
\section{Anexos}
\normalsize{ \indent
En los anexos, se presenta la declaración jurada junto con el
horario de un docente; ambos forman parte de la documentación
manejada dentro del Instituto Línea Cuchilla.
}
